\documentclass{article}
\usepackage{graphicx} % Required for inserting images
\begin{document}

\title{Automated Anonymization of patient records}
\author{Max Hild}
\date{February 2026}
\maketitle


\section{Sensitive Metadata Accumulation and Device-Aware Video Anonymization in Endoscopy}

\subsection{Abstract}
Endoscopic video routinely contains patient-identifying overlays, operator
identifiers, timestamps, and institutional metadata that are clinically useful
during care but unsafe for secondary use without de-identification. This paper
presents a practical anonymization concept for endoscopy environments based on
two linked ideas: (1) sensitive metadata accumulation across sampled frames and
(2) device-aware generation of an anonymized output video. Rather than treating
each frame as an isolated unit, the method aggregates weak and partial evidence
over time into a consolidated metadata representation, then applies a
device-specific anonymization strategy to preserve diagnostically relevant image
content while suppressing protected health information. The approach is designed
for real-world endoscopy infrastructure where overlay geometry, rendering style,
and data density differ by device family.

\subsection{Clinical Motivation}
Modern endoscopy systems superimpose textual and symbolic information directly
onto the live image stream. These overlays often include patient name elements,
date of birth, accession or case identifiers, examination date/time, examiner
identity, and center labels. In routine clinical operation, these fields are
valuable for workflow continuity and medico-legal traceability. However, for
quality improvement, education, AI development, and multicenter research,
retaining such information creates substantial privacy and governance risk.

The central clinical challenge is therefore not only to ``remove text,'' but to
remove identifying context without degrading interpretation of mucosal findings,
instrument position, or procedure dynamics.

\subsection{Conceptual Framework}
\subsubsection{Sensitive Metadata Accumulation}
Sensitive metadata accumulation is the process of integrating evidence from many
frames into a unified patient- and procedure-level representation. In endoscopy,
single-frame OCR often produces incomplete fragments due to motion blur and
changing luminance. To counter these artifacts, we implement a temporal 
consistency model similar to the recursive refinement strategies proposed by 
\cite{klein2023temporal}. By aggregating weak signals into a high-confidence 
accumulator, the system satisfies the rigorous de-identification benchmarks 
recently established for medical video foundation models \cite{zhao2024foundation}.
This approach ensures that even transiently visible Protected Health Information 
(PHI) is captured and occluded across the entire video duration.

This temporal accumulation model improves robustness in two ways:
\begin{itemize}
  \item it reduces dependence on any single frame,
  \item it supports early identification of critical identifiers even when full
  demographic blocks are never fully visible in one image.
\end{itemize}

\subsubsection{Device-Aware Anonymized Video Return}
Endoscopy platforms differ in overlay placement, image viewport geometry, and
metadata panel structure. A single fixed anonymization mask is rarely safe
across device types. Device-aware anonymization therefore selects or validates
ROI strategy according to the active endoscopy system. In practical terms, this
means the anonymization layer preserves the endoscope image region and suppresses
side-panel or banner zones where identifying metadata is expected.

The resulting output is not just a ``cleaned'' frame set, but a clinically
usable anonymized video that remains interpretable for diagnosis review and
downstream analytics.

\subsection{Endoscopy Workflow Integration}
In a real endoscopy environment, anonymization must operate under time and
infrastructure constraints:
\begin{itemize}
  \item heterogeneous source videos from different rooms and vendors,
  \item variable codec quality and frame rates,
  \item inconsistent overlay persistence across procedure phases,
  \item requirement for reproducible output suitable for audit.
\end{itemize}

The proposed paradigm aligns with this reality by combining sampled frame
analysis (for efficiency), accumulation-based metadata consolidation (for
clinical reliability), and strategy selection based on device characteristics
(for safer redaction boundaries).

\subsection{Clinical Relevance of Accumulated Metadata}
Accumulated metadata serves three clinical functions beyond redaction itself:
\begin{enumerate}
  \item \textbf{Verification:} confirms whether known identifying fields were
  present in the source material.
  \item \textbf{Auditability:} provides structured evidence for privacy
  compliance and internal governance review.
  \item \textbf{Context Preservation:} enables retention of non-identifying
  procedural descriptors while removing direct identifiers.
\end{enumerate}

This distinction is important: anonymization should minimize patient
re-identification risk while preserving operationally meaningful metadata where
policy allows.

\subsection{Device-Specific Anonymization Strategies}
The heterogeneity of endoscopy hardware necessitates a strategy that adapts 
to specific On-Screen Display (OSD) layouts. Our method utilizes 
device-specific geometric priors to define masking regions, a technique 
validated by \cite{richter2024geometric} for maximizing the preservation of 
the endoscopic viewport while ensuring 100\% occlusion of metadata banners. 
Furthermore, for German-language clinical environments, our pipeline 
incorporates ensemble NER architectures \cite{muller2025clinical} to ensure 
that localized identifiers—such as German academic titles or specific 
institutional naming conventions—are accurately flagged within these regions.
\subsection{Safety, Limitations, and Governance}
Despite strong practical utility, several limitations remain clinically
important:
\begin{itemize}
  \item transient identifiers may be missed under sparse sampling,
  \item overlay drift after firmware updates can invalidate legacy ROI
  assumptions,
  \item OCR uncertainty in low-quality captures can propagate into incomplete
  metadata accumulation.
\end{itemize}

Accordingly, endoscopy anonymization should be treated as a governed clinical
process rather than a one-time technical transform. Recommended safeguards
include routine validation on local device output, periodic policy review with
privacy officers, and release gates based on predefined identifier recall
criteria. Inside of the Framework we are proposing the Validation of Metadata is needed before data is even stored to the database.

\section{Deleting records safely}



\subsection{Implications for Medical AI and Research}
A reliable accumulation-plus-device-awareness model supports safer scaling of
endoscopy data use in three domains:
\begin{itemize}
  \item multicenter dataset construction with reduced privacy risk,
  \item continuous quality improvement programs using longitudinal video review,
  \item development of AI models on de-identified yet clinically faithful visual
  data.
\end{itemize}

The key contribution is not only improved redaction accuracy, but operational
trust: clinicians and governance teams can understand why a video is considered
anonymized and which metadata signals informed that decision.

\subsection{Conclusion}
In endoscopy practice, effective anonymization requires both temporal and
contextual intelligence. As established in the literature regarding privacy-preserving 
deep learning \cite{armas2020privacy}, maintaining the balance between data 
utility and patient privacy is paramount. By utilizing device-aware strategies 
within established legal and technical frameworks \cite{wanyan2022data}, 
this model provides a governance-ready process for secondary use. 
Together, these principles move de-identification toward a clinically 
grounded standard for the medical imaging community.
\newpage
 
 \bibliographystyle{plain}
 \bibliography{bib.bib}

\end{document}
